% ECONOMICS_PROBLEM: {"id": "D81-16", "title": "Separate identification of beliefs and preferences", "jel_code": "D81", "source_id": "OP-0587"}
% !TeX program = pdflatex
\documentclass[11pt,a4paper]{article}
\usepackage[margin=22mm]{geometry}
\usepackage[T1]{fontenc}
\usepackage[utf8]{inputenc}
\usepackage{lmodern,amsmath,amssymb,amsthm,mathtools}
\usepackage{booktabs,longtable,array,enumitem,xcolor,xurl,hyperref,listings}
\hypersetup{colorlinks=true,linkcolor=blue,urlcolor=blue}
\setlength{\parindent}{0pt}
\setlength{\parskip}{5pt}
\setlength{\emergencystretch}{3em}
\newtheorem{theorem}{Theorem}[section]
\newtheorem{lemma}[theorem]{Lemma}
\newtheorem{proposition}[theorem]{Proposition}
\newtheorem{corollary}[theorem]{Corollary}
\theoremstyle{definition}\newtheorem{definition}[theorem]{Definition}
\theoremstyle{remark}\newtheorem{remark}[theorem]{Remark}
\newcommand{\R}{\mathbb R}\newcommand{\Q}{\mathbb Q}
\newcommand{\N}{\mathbb N}\newcommand{\E}{\mathbb E}
\newcommand{\1}{\mathbf 1}
\DeclareMathOperator*{\argmax}{arg\,max}
\DeclareMathOperator*{\argmin}{arg\,min}
\newcommand{\supp}{\operatorname{supp}}
\newcommand{\conv}{\operatorname{conv}}
\newcommand{\rank}{\operatorname{rank}}
\lstset{basicstyle=\ttfamily\scriptsize,breaklines=true,columns=fullflexible,
 keepspaces=true,showstringspaces=false,frame=single,aboveskip=8pt,belowskip=8pt}
\begin{document}
\begin{center}
{\Large\bfseries Separate identification of beliefs and preferences}\par
\medskip
{\large D81-16: research attempt and adversarial verification}\par
9 October 2026
\end{center}
\textbf{Tools.} Web browsing: YES. Exact rational computation: YES.\par
\section{FORMAL RESTATEMENT}
Let $\Omega=[m]$, $m\ge1$, $p\in\Delta_m=\{p\in\R_+^m:\sum_sp_s=1\}$, and let $\{u_\theta:\theta\in\Theta\}$ be a specified increasing utility family on consequences $C$. Fix $c_0<c_1$ in $C$ and require $u_\theta(c_0)=0$, $u_\theta(c_1)=1$. A supported environment $z\in Z$ provides a nonempty finite menu $A(z)$ and known consequences $c(a,s,z)$. Expected utility is
\[
 U_{p,\theta}(a,z)=\sum_{s=1}^m p_su_\theta(c(a,s,z)).
\]
A declared observation mechanism induces a map $\mathcal O$ from $(p,\theta)$ to a population law of observed choices and elicited reports. Identification of the pair means
\[
 \forall(p,\theta),(p',\theta'):\quad
 \mathcal O(p,\theta)=\mathcal O(p',\theta')
 \Longrightarrow (p,\theta)=(p',\theta').
\]
Identification of a functional $\psi$ replaces the conclusion by equality of $\psi$.

\textbf{Literal statement versus a specified theorem.} The attachment asks for a characterization across unspecified utility families, menu designs, elicitation mechanisms, and ambiguity criteria. It is a research agenda, not an assertion that all such designs identify the pair. There is no single fixed $\mathcal O$ whose injectivity can be decided without specifying these primitives. A nonidentifying example therefore does not disprove the agenda. The minimal operational specialization used below observes the \emph{full deterministic maximizing correspondence} at each supported menu, without choice error. Additional exact reports, when present, have a known map $E(p,\theta)$. The lottery and ambiguity theorems declare further assumptions rather than importing them silently.

\textbf{Stress points.} Positive-affine normalization does not remove duplicate parameter labels that represent the same normalized utility. Finite strict choices supply inequalities, not indifference equations. Observing a selected maximizer is weaker than observing the whole maximizing set. An elicited probability is not assumed truthful merely because a survey labels it a belief. Priors may have zero coordinates. Objective lottery evaluation and invariance of beliefs/preferences across interventions are substantive restrictions.

\section{QUICK LITERATURE/CONTEXT CHECK}
Charles F.~Manski, \emph{Survey Measurement of Probabilistic Macroeconomic Expectations: Progress and Promise}, \emph{NBER Macroeconomics Annual} 32 (2018), 411--471, discusses choice without expectations data in Section II.A and confounding beliefs and preferences in Section VII.C. Checked publisher text: \url{https://www.journals.uchicago.edu/doi/10.1086/696061}. The latter section also explains why expectations reports require care when beliefs and preferences are not separately represented. It does not state the attachment's map-injectivity question as a numbered theorem or conjecture.

The results below are mathematical statements for explicitly specified submodels. The rich lottery model uses expected utility for objective randomization as an assumption, not as an empirical consequence of the observed choices. No general identification result is attributed to the cited discussion.

\section{ATTACK PLAN}
\textbf{Proof track.} Write the sharp observational fibers exactly; design utility-calibrating comparisons independent of beliefs and belief-calibrating comparisons with normalized utilities; identify which ambiguity-set object is separated by all such comparisons.

\textbf{Disproof track.} Use finitely many strict choices to test local identification; construct two different normalized utility/prior pairs with the same entire continuum of certainty comparisons; compare distinct prior sets with the same lower expectation functional.

\textbf{Tools and roles.} Revealed-preference inequalities describe sharp sets; finiteness and continuity create local observational equivalence; explicit power utilities exhibit confounding; normalized binary lotteries calibrate a common cardinal scale; exclusion restrictions hold latent parameters fixed; support-function separation identifies closed convex prior sets; compactness guarantees a separating nearest point; fiber factorization characterizes identified functionals.

\textbf{Tiny cases.} With $m=1$, the prior is fixed at one but utility parameters need not be identified. Menus containing only the reference outcomes reveal no curvature parameter. With two states, even the continuum of menus in Proposition~\ref{prop:cutoff} fails to identify either primitive separately.

\section{WORK}
\begin{theorem}[Sharp deterministic identified set]\label{thm:fiber}
Suppose the full nonempty maximizing set $B(z)\subseteq A(z)$ is observed for every $z$, and an exact report $e$ has known structural map $E$. Then the sharp set of compatible parameters consists exactly of $(p,\theta)\in\Delta_m\times\Theta$ satisfying $E(p,\theta)=e$ and, for every $z$,
\[
 \begin{aligned}
 U_{p,\theta}(a,z)&=U_{p,\theta}(a',z)
       &&(a,a'\in B(z)),\\
 U_{p,\theta}(a,z)&>U_{p,\theta}(b,z)
       &&(a\in B(z),\ b\in A(z)\setminus B(z)).
 \end{aligned}
\]
If no report is observed, omit its equality. A functional is point identified at these data if and only if it is constant on this set; its sharp range is its image of this set.
\end{theorem}
\begin{proof}
Every maximizing action has the same maximum utility. An action outside the full maximizing set has strictly smaller utility, because the menu is finite and includes a maximizer. This proves necessity. Conversely, the displayed equalities and strict inequalities make the full maximizing set exactly $B(z)$ at every menu. Together with the report equation, this reproduces every component of the declared observation. This proves sufficiency and sharpness. Two feasible parameter pairs giving different functional values show nonidentification of that functional. If every feasible pair gives the same value, that value is determined by the data. Taking all feasible functional values gives exactly the compatible range.
\end{proof}

\textbf{Scope of the inequalities.} If only one selected action is observed and the tie rule is unobserved, weak inequalities rationalize the selection but do not recover the full-correspondence identified set above. A known tie rule or error distribution must instead enter the observation map explicitly. Similarly, imposing $p=e$ on a reported vector is valid only when truthful, exact elicitation is an assumption of the model.

\begin{proposition}[Finite strict choices cannot locally identify continuous parameters]\label{prop:finite}
Let $Z$ be finite and all $A(z)$ finite. Suppose all pairwise utility differences are continuous in a neighborhood of a nonisolated parameter point $(p^0,\theta^0)$, and every observed menu has a unique maximizer there. If no additional reports distinguish nearby parameters, there is a relative neighborhood of $(p^0,\theta^0)$ with the same choices. Hence the point is not identified.
\end{proposition}
\begin{proof}
For each observed maximizer and each competing action, its utility advantage at the specified point is strictly positive. There are finitely many such differences. Continuity gives a neighborhood on which each remains positive; their finite intersection is a neighborhood on which every maximizing action is unchanged. By nonisolation it contains a distinct admissible parameter point. In the absence of distinguishing reports that point gives identical observations. If a menu is a singleton, it contributes no comparisons and imposes no restriction.
\end{proof}
Without continuity, there is still a global cardinal obstruction for a finite deterministic choice-only design with no additional reports: each finite menu has finitely many possible nonempty maximizing subsets, so the combined observation alphabet is finite. It cannot be the injective image of an infinite parameter set. This does not rule out identification at exceptional points where indifferences are observed, nor does it apply to a continuous-valued observation law generated by a specified stochastic choice model.

\begin{proposition}[A continuum of menus can still confound beliefs and utility]\label{prop:cutoff}
Let $m=2$, $C=[0,1]$, $u_\theta(c)=c^\theta$ with $\theta>0$, and prior $(p,1-p)$ with $0<p<1$. For every $c\in[0,1]$, observe the full choice correspondence between the bet $B=(1,0)$ and the sure consequence $c$. This entire continuum of observations identifies exactly
\[
 t=p^{1/\theta}\in(0,1),
 \qquad\{(p,\theta):p=t^\theta,\ \theta>0\}
\]
as the observational fiber. In particular $(p,\theta)=(1/2,1)$ and $(1/4,2)$ are observationally equivalent.
\end{proposition}
\begin{proof}
The bet has expected utility $p$ and the sure consequence has utility $c^\theta$. Strict monotonicity of $c^\theta$ implies that the bet is strictly preferred for $c<t$, both are chosen at $c=t$, and the sure consequence is strictly preferred for $c>t$. Hence $t$ determines every observation. Conversely, different cutoffs yield different choices at any consequence strictly between them, so the observation determines $t$. Solving $t=p^{1/\theta}$ gives the displayed exact fiber. Both stated pairs have cutoff $1/2$, and each utility satisfies the required normalization. As $\theta$ ranges over $(0,\infty)$ with fixed interior $t$, $t^\theta$ ranges over $(0,1)$, so the prior itself is unrestricted within that open interval by this design.
\end{proof}
Within this submodel, a functional is identified precisely when it is constant along every relevant cutoff fiber. The cutoff itself is identified; $p$ and $\theta$ are not. This is a counterexample to an unrestricted identification claim, not to the request to characterize informative designs.

\begin{theorem}[Separate identification under a declared rich lottery design]\label{thm:lottery}
Assume $C=[c_0,c_1]$. Objective lotteries are evaluated by the same normalized von Neumann--Morgenstern utility as sure consequences, with their stated probabilities; randomization does not change $p$ or $u_\theta$. Let $L_\lambda$ pay $c_1$ with objective probability $\lambda\in[0,1]$ and $c_0$ otherwise, independently of the state. Suppose the following two-action choice correspondences are observed for all indicated parameters:

(a) every sure consequence $c\in C$ versus $L_\lambda$, for every $\lambda\in[0,1]$;

(b) for each state $s$, the bet $B_s$ paying $c_1$ in state $s$ and $c_0$ otherwise, versus every $L_\lambda$.

Then $p$ and the normalized function $u_\theta$ on $C$ are identified separately, including priors on the boundary of the simplex. The label $\theta$ is identified if and only if the utility family has no observationally duplicate labels, that is, $\theta\mapsto u_\theta|_C$ is injective on the allowed parameter set.
\end{theorem}
\begin{proof}
Normalization and expected utility make $U(L_\lambda)=\lambda$, irrespective of $p$. Monotonicity and the endpoint normalization give $u_\theta(c)\in[0,1]$. In menu (a), the sure consequence is weakly preferred exactly when $u_\theta(c)\ge\lambda$. The full set of such $\lambda$, including its endpoints, is $[0,u_\theta(c)]$, so its supremum recovers $u_\theta(c)$. This works for each $c$ without any continuity or differentiability of $u_\theta$.

The state bet has utility $\sum_t p_t\1\{t=s\}=p_s$. Menu (b) therefore gives the interval $[0,p_s]$ of lottery probabilities for which the bet is weakly preferred. Its supremum recovers $p_s$, including $p_s=0$ or $1$. The same objective lotteries supply a common numerical scale for both steps.

If the family is injectively labeled, the recovered utility function determines $\theta$. If two labels give the same utility function, then for any fixed $p$ every menu in this design has the same utility comparisons under both labels; these labels cannot be distinguished. This proves the stated necessity as well as sufficiency.
\end{proof}

\textbf{Exclusion restriction made explicit.} The theorem holds $p$ and the entire utility function fixed across the lottery probabilities, consequence values, and state-bet menus. The interpretation fails if changing a menu changes subjective states, beliefs, preferences, or evaluation of the objective probabilities. No choice observation in the theorem independently proves that invariance.

\begin{theorem}[What max--min ambiguity data identify]\label{thm:ambiguity}
Let $\varnothing\ne\Pi\subseteq\Delta_m$ and declare the criterion
\[
 V_{\Pi,\theta}(f)=\inf_{p\in\Pi}\sum_s p_s\E_{f(s)}u_\theta(c),
\]
where each $f(s)$ is a known objective lottery on $C=[c_0,c_1]$. Keep the lottery evaluation and invariance assumptions of Theorem~\ref{thm:lottery}. Observe its utility-calibrating menus and, for every $x\in[0,1]^m$, comparisons of the act $f_x(s)=L_{x_s}$ with every constant lottery $L_\lambda$.

These observations identify $u_\theta$ and exactly the closed convex set
$K=\overline{\conv}(\Pi)$. If $\Pi$ is required to be closed and convex, it is identified itself. Without that representation restriction, observationally equivalent prior sets are exactly those with the same closed convex hull, once the normalized utility is fixed.
\end{theorem}
\begin{proof}
For state-independent consequences and lotteries, taking an infimum over $p$ leaves their utility unchanged since $\sum_sp_s=1$. The utility-calibration proof therefore still applies. For $f_x$ the value is
\[
 \underline E_\Pi(x)=\inf_{p\in\Pi}p\cdot x\in[0,1].
\]
Comparisons with $L_\lambda$ identify this value as the supremum of lottery probabilities weakly below it.

A linear functional has the same infimum over a set and over its convex hull: each convex combination has a value at least the original infimum, and the original set is contained in the hull. Continuity preserves this infimum on taking the closure. Because the simplex is compact, $K$ is nonempty compact convex and
$\underline E_\Pi(x)=\min_{p\in K}p\cdot x$.

It remains to show that all the recovered lower expectations determine $K$. First they determine $\min_{p\in K}p\cdot v$ for every $v\in\R^m$. If $v$ is constant its value follows from $\sum p_s=1$. Otherwise put $a=\min_sv_s$, $b=\max_sv_s-a>0$, and $x=(v-a\1)/b\in[0,1]^m$. Then
\[
 \min_{p\in K}p\cdot v=a+b\min_{p\in K}p\cdot x.
\]
Now suppose two nonempty compact convex sets $K,K'\subseteq\Delta_m$ differ. After interchanging their names if needed, choose $p^0\in K\setminus K'$, and choose $q^0\in K'$ minimizing the Euclidean distance to $p^0$. Compactness gives the minimizer and the distance is positive. Let $h=p^0-q^0$. For $q\in K'$, convexity puts $q^0+t(q-q^0)$ in $K'$ for $0\le t\le1$. Minimality of $q^0$ and the derivative at $t=0$ of its squared distance to $p^0$ yield $h\cdot(q-q^0)\le0$. Thus
\[
 h\cdot p^0=h\cdot q^0+\|h\|^2>\max_{q\in K'}h\cdot q.
\]
Taking $v=-h$ gives
$\min_{p\in K}p\cdot v\le p^0\cdot v<\min_{q\in K'}q\cdot v$,
contradicting equality of all recovered functionals. Hence $K$ is unique. Conversely, equal closed convex hulls give equal infima for every state utility vector and therefore identical choices for all acts under this declared criterion. The remaining identification of a utility parameter label has the same injectivity condition as before.
\end{proof}

\textbf{Explicit ambiguity-set nonidentification.} For two states, $\Pi=\{(1,0),(0,1)\}$ and $\Pi'=\Delta_2$ are distinct, but both give $\min(x_1,x_2)$ for every state utility vector $x$. No amount of these max--min choice data can distinguish their raw representations. Their common closed convex hull is identified by Theorem~\ref{thm:ambiguity}. The theorem does not apply without modification to smooth ambiguity, variational penalties, or other undeclared criteria.

\section{VERIFICATION}
\textbf{Exact tests.} The cutoff example was checked at the 101 rational consequences $c=k/100$, retaining ties. The same-sign identity also follows algebraically from
$1/4-c^2=(1/2-c)(1/2+c)$, whose second factor is strictly positive on $[0,1]$.
\begin{lstlisting}[language=Python]
from fractions import Fraction as F
sgn=lambda x:(x>0)-(x<0)
for k in range(101):
    c=F(k,100)
    assert sgn(F(1,2)-c)==sgn(F(1,4)-c*c)
for x1 in range(6):
    for x2 in range(6):
        a,b=F(x1,5),F(x2,5)
        # Any mixture of the two state values lies above their minimum.
        for k in range(11):
            p=F(k,10)
            assert p*a+(1-p)*b>=min(a,b)
print('101 cutoff comparisons and ambiguity-mixture grid verified exactly')
\end{lstlisting}
The continuum results do not rely on extrapolating these finite tests; their proofs establish the comparisons for every admissible consequence and probability.

\textbf{Adversarial audit.} The sharp set uses strict inequalities only because the entire maximizing correspondence is observed. The finite-choice local theorem explicitly excludes extra continuous reports that might identify parameters. The positive-affine normalization is shared across menus, and the utility calibration uses state-independent objective lotteries, not an unproved identification of their probabilities. The prior-set theorem identifies a closed convex hull, not an arbitrary generating set. All compactness is finite dimensional. No statistical consistency, confidence interval, or sampling-rate claim is inferred from population identification.

\section{FINAL}
\textbf{LABEL: UNRESOLVED for the full research agenda as stated.}

\textbf{Strongest fully proved partial result.} Exact sharp sets are derived for deterministic full-correspondence observations. Separate identification is proved for the specified rich lottery design, and closed-convex-hull identification is proved for the declared max--min ambiguity criterion. Explicit normalized observationally equivalent pairs disprove unrestricted separate identification even with a continuum of sure-consequence comparisons. These are complete submodel theorems, not a characterization of every unspecified observation mechanism.

\textbf{Exact first gap.} No necessary-and-sufficient characterization in primitive restrictions is proved for an arbitrary supported-menu design, utility family, elicitation/error mechanism, and ambiguity criterion. Outside the specified submodels these inputs remain free, so the tautological requirement that observations separate parameters has not been converted into such a characterization.

\textbf{Three next targets.} First, fix a finite-dimensional utility family and a concrete menu design and characterize when observed indifference equations have a unique global solution, not merely a full-rank local derivative. Second, specify an elicitation measurement law and derive the joint choice/report fiber rather than assuming reports equal true priors. Third, repeat the lower-functional separation analysis for one declared non-max--min ambiguity criterion, identifying its representation equivalences before attempting parameter recovery.

\textbf{Minimal nonidentification structure.} Two states and a one-dimensional normalized utility family already suffice, with positive priors and no choice error. Even one state can leave utility unidentified under poor menu variation. A candidate broader theorem must exclude these designs or identify only functionals constant on their displayed fibers.

\end{document}
